\section{Compute resources}
\label{app:compute}

All experiments ran on a cloud virtual machine with two CPU cores and 7\,GB of memory, without a GPU. Feature extraction for one release takes minutes; one run of the structure view over $13{,}600$ images takes a few minutes; each fine-tuning run took $4$--$9$ minutes of training ($143$ minutes for all $24$ runs). Exploratory runs that are not reported (for example an earlier three-class CIFAR-Aug subset, Appendix~\ref{app:hyper}) used comparable resources.

\section{Licences of existing assets}
\label{app:licences}

CIFAR-10~\citep{krizhevsky2009cifar}: no licence is stated by its creators; it is distributed for research use. PlantDoc~\citep{singh2020plantdoc}: CC BY 4.0. COVID-19 image data collection~\citep{cohen2020covid}: a licence is given per image in its metadata (mostly CC BY, CC BY-NC-SA and CC BY-NC-ND). Brain tumour MRI~\citep{bhuvaji2020brain}: permissive terms set by its creators (CC0 on Kaggle, MIT on the creators' Hugging Face copy). Fruits-360~\citep{muresan2018fruits}: MIT. The two leaf releases are distributed under the licences stated on their public repositories, which we will name after review. ImageNet weights come from timm~\citep{wightman2019timm} (Apache 2.0).

\section{Use of large language models}
\label{app:llm}

Generative AI tools were used extensively: to help design the experiments, to write and debug the analysis code, to draft and edit the text and to produce the figures. They also helped search the related literature; every reference was then checked against its original source. The visual inspection of candidate duplicate pairs (Section~\ref{sec:audit}) was performed with AI assistance; every inspected pair and judgement is released so that it can be checked. All numbers in the paper are produced by the released code from the released data. The authors reviewed all AI-assisted work and take full responsibility for the content.
